% Garbage collector. Requires tikz_styles.tex.
\begin{figure}[htbp]
\centering
\begin{tikzpicture}[node distance=2.6mm and 3mm]
  \node[trap, minimum width=42mm] (fail) {allocation does not fit\\abort before commit};
  \node[acc, right=of fail] (alloc) {S\_GC\_ALLOC};
  \node[acc, below=of alloc] (enter) {S\_GC\_ENTER\\drain the dmem beat};
  \node[acc, below=of enter] (roots) {S\_GC\_ROOTS\\registers and RF};
  \node[acc, below=of roots] (run) {S\_GC\_RUN\\mark and sweep};
  \node[rf, below=of run] (resume) {pop a free run\\re-dispatch};
  \draw[arr] (fail) -- (alloc);
  \draw[arr] (alloc) -- (enter);
  \draw[arr] (enter) -- (roots);
  \draw[arr] (roots) -- (run);
  \draw[arr] (run) -- (resume);
  \node[anno, align=left, anchor=north west] at ($(fail.south west)+(0,-2.6)$)
    {Off unless \texttt{+GC\_EN=1}.\\
     With the collector off, an empty heap\\
     region is still the bump pointer.};
\end{tikzpicture}
\caption{Where a collection sits in the core.
Collections start at instruction boundaries.
An allocation that does not fit aborts that instruction before the first architectural commit, the core collects, and the instruction is fetched again.
\texttt{S\_GC\_ENTER} waits until any outstanding data beat from the abort has acknowledged, so a skip-clear cannot treat that read data as the first word of the static prune map.
If the free-run list still has a run, \texttt{S\_GC\_ALLOC} satisfies the request without a collection.}
\label{fig:gc-core}
\end{figure}

\begin{figure}[p]
\centering
\begin{tikzpicture}[node distance=3mm and 4mm]
  \node[ctl] (roots) {root stream};
  \node[acc, right=of roots] (tr) {tracer};
  \node[rf, right=of tr] (q) {queue\\4 entries};
  \node[acc, right=of q] (mk) {marker};
  \node[mem, below=of mk] (bm) {on-chip bitmap\\one 128-bit word / cycle};
  \node[rf, left=of bm] (stk) {mark stack\\256 on chip};
  \node[mem, left=of stk] (spill) {spill 32\\to dmem};
  \node[mem, below=of spill] (re) {rescan list\\4096};
  \draw[arr] (roots) -- (tr);
  \draw[arr] (tr) -- (q);
  \draw[arr] (q) -- (mk);
  \draw[arr] (mk) -- (bm);
  \draw[arr] (mk) -- (stk);
  \draw[arr] (stk) -- (spill);
  \draw[arr] (mk) -- (re);
  \node[anno, align=left, anchor=north] at ($(re.south)+(0,-0.15)$)
    {One outstanding request. The marker wins a tie with the tracer.};
\end{tikzpicture}
\caption{Mark engine inside \texttt{pycore\_gc}.
The tracer walks memory-resident roots and the slots of each popped object, and emits child handles into the queue.
The marker tests the first granule of an extent, sets every granule, and pushes kinds that have children.
The bitmap is the on-chip array \texttt{bm\_q} (\texttt{PYCORE\_GC\_BITMAP\_WORDS} words, one bit per 16-byte granule).
The 256-entry on-chip stack spills its oldest 32 entries to \texttt{0xFC0000} when full and refills 32 when empty.
A pop scans at most 128 slots; the rest of the range is pushed first as a continuation, so a wide container holds at most 129 stack entries.
A child that still does not fit is left unmarked and the range is written once to the rescan list.
When the stack drains, that range is scanned again.
The collection is abandoned only when the rescan list itself is full.}
\label{fig:gc-mark}
\end{figure}

\begin{figure}[htbp]
\centering
\begin{tikzpicture}[node distance=2.4mm]
  \node[mem, minimum width=68mm] (sw) {sweep bitmap words up to the high-water word};
  \node[ctl, minimum width=68mm, below=of sw] (skip) {words above the high-water mark are one free run};
  \node[rf, minimum width=68mm, below=of skip] (keep) {jump the kept run; do not paint it};
  \node[acc, minimum width=68mm, below=of keep] (clr) {write zero back to every non-zero word};
  \node[mem, minimum width=68mm, below=of clr] (hdr) {free-run header \quad first 1024 on chip};
  \node[mem, minimum width=68mm, below=of hdr] (ov) {overflow headers in the run table, then in place};
  \foreach \p/\q in {sw/skip,skip/keep,keep/clr,clr/hdr,hdr/ov} \draw[arr] (\p) -- (\q);
\end{tikzpicture}
\caption{Sweep and the free-run list.
The sweep clears the bitmap as it scans, so the next collection starts from a clean map.
After reset, and after an abandoned collection, the idle engine clears the bitmap in the background, one word per cycle.
Free runs are not a moving collector: a live object stays at its address.
The first 1024 run headers stay on chip.
Further headers go to \texttt{PYCORE\_GC\_RUN\_TABLE}.
Past that table they occupy the run's first granule, so the list has no fixed cap.
Allocator traffic in \texttt{S\_GC\_ALLOC} reads a header and zero-fills the granted line.}
\label{fig:gc-sweep}
\end{figure}

\begin{figure}[htbp]
\centering
\begin{tikzpicture}[node distance=1.7mm and 2mm]
  \node[mem, font=\tiny\sffamily, minimum width=22mm] (k0) {K\_TUPLE};
  \node[mem, font=\tiny\sffamily, minimum width=22mm, right=of k0] (k1) {K\_CODE};
  \node[mem, font=\tiny\sffamily, minimum width=22mm, right=of k1] (k2) {K\_LIST};
  \node[mem, font=\tiny\sffamily, minimum width=22mm, right=of k2] (k3) {K\_DICT};
  \node[mem, font=\tiny\sffamily, minimum width=22mm, below=of k0] (k4) {K\_SET};
  \node[ink, font=\tiny\sffamily, minimum width=22mm, right=of k4] (k5) {K\_OBJ};
  \node[mem, font=\tiny\sffamily, minimum width=22mm, right=of k5] (k6) {K\_STR};
  \node[acc, font=\tiny\sffamily, minimum width=22mm, right=of k6] (k7) {K\_DICTT};
\end{tikzpicture}
\caption{Mark-stack entry kinds.
An entry is 67 bits: kind, size, and address.
\texttt{K\_TUPLE} is also the continuation kind for 32-byte slots (tuples, list buffers, set tables, dict order).
\texttt{K\_DICTT} continues a dict hash table, whose slots are 64 bytes.
\texttt{K\_OBJ} covers the \texttt{OBJECT} header kinds: instance, type, bound method, builtin, bytearray, exception, cell, and function.
Strings are traced for the object header; the payload bytes hold no pointers.}
\label{fig:gc-kinds}
\end{figure}
